```latex
\documentclass{article}
\usepackage[spanish]{babel}
\usepackage[utf8]{inputenc}
\usepackage{amsmath}
\usepackage{hyperref}
\usepackage{listings}
\usepackage[margin=1in]{geometry}

\hypersetup{
    colorlinks=true,
    urlcolor=blue,
    linkcolor=blue,
    citecolor=blue
}

\lstset{
    language=[ANSI]C,
    basicstyle=\ttfamily\small,
    keywordstyle=\color{blue},
    commentstyle=\color{gray},
    stringstyle=\color{red},
    breaklines=true,
    frame=single,
    captionpos=b,
    columns=fullflexible,
    keepspaces=true
}

\makeatletter
\def\lstlistingname{Código}
\makeatother

\title{Apuntes de la Asignatura JK: Tema 1 --- Arquitectures i Eines}
\author{Profesor Experto en Informática}
\date{\today}

\begin{document}
\maketitle

\section{Introducció al disseny d'arquestures computing}
El Tema 1 de la assignatura JK s'adentra en els conceptes fonamentals que subyacen al disseny i la comprensió d'arquestures de computadors modernes. Aquest capítol proporciona una base rigurosa per a comprendre com es organitzen els sistemes informàtics des del nivell de hardware fins a les eines de desenvolupament software.

\subsection{Arquestura de Von Neumann i Princeton}
La arquestura de Von Neumann, proposta l'any 1945, és el model bàsic sobre el qual es construeixen la majoria dels ordinadors generals d'actualdia. El seu principi és l'emmagatzematge conjoint de dades i instruccions en una memòria unificada, accessiblevia mitjançant una adreça comú.

\begin{itemize}
    \item \textbf{Memòria unificada:} Tant les instruccions com els dades residen en la mateixa memòria RAM.
    \item \textbf{Adreçament:} El processador utilitza un contador de programa (Program Counter, PC) per a obtenir la següent instrucció des de la memòria.
    \item \textbf{Bus de von Neumann:} El mateix bus s'utilitza per a transferir instruccions i dades, el que pot crear un "coll de botella" (von Neumann bottleneck).
\end{itemize}

En contraposició, l'arquestura de Princeton separa físicament la memòria d'instruccions i la memòria de dades, permetent accesos simultànies (fetch i read/write), però és menys flexible en termes de gestió dinàmica de memòria.

\subsection{Components principals d'un sistema}
Tots els sistemes computing, independentment de la seva mida, comparten un conjunt de components bàsics:

\begin{center}
\begin{tabular}{lll}
\textbf{Component} & \textbf{Funció} & \textbf{Detall} \\
CPU (Central Processing Unit) & Executar instruccions & ALU i Unitat de Control \\
Memòria Principal (RAM) & Almacenament temporal & Accés aleatori, volatil \\
Memòria Secundària (HDD/SSD) & Almacenament persistent & No volàtil, major capacitat \\
Units d'E/S (Input/Output) & Interfície amb el món exterior & Teclat, ratolí, xarxa, gràfiques \\
\end{tabular}
\end{center}

\subsection{Conjunts d'instruccions (ISA)}
L'ISA defineix el vocablar màxim que un processador entén. Inclou:
\begin{itemize}
    \item Tipus de dades (enter senzill, enter senzill no signat, de pla, etc.).
    \item Formats d'instrucció (R-type, I-type, S-type en MIPS).
    \item Registres visibles per al programador.
    \item Sistemes d'adreçament.
\end{itemize}

\subsection{Eines de desenvolupament}
Per a treballar ambarquitectures de baix nivell, disposem d'un conjunt d'eines:
\begin{itemize}
    \item \textbf{Assemblers:} Traduixen llenguatge ensamblador a codi màquina. Exemples: \texttt{nasm}, \texttt{mips-as}.
    \item \textbf{Simuladors d'arquitectura:} Permeten executar codi màquina en un entorn controlat per estudiar comportament. Exemple: \texttt{simplesim}, \texttt{gem5}.
    \item \textbf{Debuggers:} Faciliten la inspecció i modificació del flux d'execució pas a pas. Exemple: \texttt{gdb}.
    \item \textbf{Compiladors:} Traduixen llenguatges de alt nivell a codi màquina o ensamb